% Author: Writer agent · 2026-11-14
% Task W1: Chapter 3 (collaborative filtering + cosine lab). Sources: research-notes.md, Ch. 3.
% Lab numbers (FICTIONAL): Noa (5,2,1), Yoni (4,3,0), Maya (1,0,5).
% Noa·Yoni = 26, |Noa| = sqrt(30) ≈ 5.48, |Yoni| = 5, cos = 26/27.39 ≈ 0.95.
% Noa·Maya = 10, |Maya| = sqrt(26) ≈ 5.10, cos = 10/sqrt(780) = 10/27.93 ≈ 0.36.

\chapter{אנשים כמוך}\label{ch:3}

\begin{scene}
יום שישי בערב, בר קטן בפלורנטין. יוני מניח את הטלפון על השולחן בחגיגיות: „האפליקציה המליצה לי אתמול על שיר שהוא פשוט את, נועה. פופ, דרמה, קלידים מוגזמים".

„גם לי!" אומרת מאיה. „ואני אפילו לא אוהבת אותו".

נועה מביטה בשניהם בחשד. „רגע. האפליקציה מרכלת עליי?"

„לא", אומר יוני. „היא פשוט יודעת שאנחנו דומים".

„אנחנו לא דומים", אומרת נועה, ומזמינה בדיוק את מה שיוני הזמין.
\end{scene}

\section{מי שהסכימו בעבר}

הרעיון של „אנשים כמוך" ותיק יותר מאפליקציות המוזיקה. כבר ב-1992 הופיע המונח \term{סינון שיתופי}{collaborative filtering} בכותרת של מאמר שתיאר מערכת בשם \en{Tapestry} \cite{goldberg1992using}. ב-1994 תיארו חוקרי \en{GroupLens} מערכת שמנבאת ציונים לפי הציונים שנתנו אנשים אחרים. לדבריהם, מסננים שיתופיים עוזרים לאנשים לבחור על סמך דעותיהם של אחרים, מתוך הנחה פשוטה: מי שהסכימו בעבר כנראה יסכימו גם בעתיד \cite{resnick1994grouplens}.

סינון שיתופי משתמש בדמיון בין משתמשים ובין פריטים בו בזמן \cite{google2026candidate}. המערכת לא צריכה לדעת מה זה „פופ" או מה זה „קלידים מוגזמים". מספיק לה לדעת מי שמע מה.

\section{טבלה ענקית של מי שמע מה}

כדי לדמיין את זה, חשבו על טבלה ענקית: כל שורה היא מאזין, כל עמודה היא שיר, ובכל תא כתוב כמה פעמים המאזין שמע את השיר. ברור שרוב התאים ריקים, כי אף אחד לא שמע 100 מיליון שירים. כל שורה היא רשימה של מספרים, ובמתמטיקה קוראים לרשימה כזו \term{וקטור}{vector}. כדי למצוא „אנשים כמוך" צריך דרך למדוד עד כמה שני וקטורים דומים. אחת הדרכים המוכרות היא \term{דמיון קוסינוס}{cosine similarity}.

\begin{lab}{דמיון קוסינוס}
זו נוסחה לימודית, ולא האלגוריתם של חברה כלשהי. דמיון קוסינוס הוא מדד דמיון סטנדרטי: „פשוט הקוסינוס של הזווית בין שני הווקטורים". יש גם מדדים אחרים, כמו המכפלה הסקלרית והמרחק האוקלידי \cite{google2026candidate}.
\[\cos(\theta)=\frac{\mathbf{a}\cdot\mathbf{b}}{\lVert\mathbf{a}\rVert\,\lVert\mathbf{b}\rVert}\]
\textbf{מה כל סימן אומר:}
\begin{itemize}
  \item $\mathbf{a}$ ו-$\mathbf{b}$: שני הווקטורים שמשווים. כאן, השורות של שני מאזינים בטבלה.
  \item $\mathbf{a}\cdot\mathbf{b}$: המכפלה הסקלרית. כופלים את המספרים שבאותו מקום (שיר מול שיר) ומחברים את כל המכפלות.
  \item $\lVert\mathbf{a}\rVert$ ו-$\lVert\mathbf{b}\rVert$: האורך של כל וקטור. מעלים כל מספר בריבוע, מחברים, ומוציאים שורש.
  \item $\theta$: הזווית בין שני הווקטורים. ככל שהיא קטנה יותר, הטעמים „מצביעים" לאותו כיוון.
  \item $\cos(\theta)$: התוצאה. כשכל המספרים אינם שליליים, כמו מספרי השמעות, היא בין 0 ל-1. ערך 1 פירושו טעם זהה ביחסים, וערך 0 פירושו שאין שום שיר משותף.
\end{itemize}

% Layout via class commands \besidetable and \sidebyside (Prompt 18; audit A2-13).
\besidetable{%
\begin{tabular}{@{}lccc@{}}
\toprule
מאזין/ה & שיר א & שיר ב & שיר ג\\
\midrule
נועה ($\mathbf{n}$) & 5 & 2 & 1\\
יוני ($\mathbf{y}$) & 4 & 3 & 0\\
מאיה ($\mathbf{m}$) & 1 & 0 & 5\\
\bottomrule
\end{tabular}}{%
\textbf{דוגמה מחושבת.} הטבלה שמימין מראה כמה פעמים כל אחד שמע שלושה שירים בשבוע אחד, והיא בדיונית לגמרי \fictional.

האורך של וקטור ההשמעות של נועה משותף לשתי ההשוואות:
$\lVert\mathbf{n}\rVert=\sqrt{5^2+2^2+1^2}=\sqrt{30}\approx5.48$.}

\sidebyside{%
\textbf{נועה מול יוני.}
\[\mathbf{n}\cdot\mathbf{y}=5\cdot4+2\cdot3+1\cdot0=26\]
\[\lVert\mathbf{y}\rVert=\sqrt{4^2+3^2+0^2}=\sqrt{25}=5\]
\[\cos(\theta)=\frac{26}{\sqrt{30}\cdot5}\approx\frac{26}{27.39}\approx0.95\]}{%
\textbf{נועה מול מאיה.}
\[\mathbf{n}\cdot\mathbf{m}=5\cdot1+2\cdot0+1\cdot5=10\]
\[\lVert\mathbf{m}\rVert=\sqrt{1^2+0^2+5^2}=\sqrt{26}\approx5.10\]
\[\cos(\theta)=\frac{10}{\sqrt{30}\cdot\sqrt{26}}\approx\frac{10}{27.93}\approx0.36\]}

\textbf{המסקנה.} יוני (0.95) דומה לנועה הרבה יותר ממאיה (0.36). שניהם שומעים בעיקר את שיר א, ומאיה שומעת בעיקר את שיר ג. מערכת שעובדת בשיטת „אנשים כמוך" תמליץ לנועה קודם כול על שירים שיוני אוהב ונועה עוד לא שמעה. שימו לב: הקוסינוס בודק את הכיוון ולא את הכמות. אילו יוני היה שומע כל שיר פי שניים, התוצאה הייתה נשארת 0.95.
\end{lab}

\section{משכנים לגורמים נסתרים}

השוואה ישירה בין מאזינים, כמו בדוגמה, היא הרעיון שמאחורי שיטות „השכנים הקרובים". לפי קורן, בל ווולינסקי, תחרות \en{Netflix Prize} הראתה שמודלים של \term{פירוק מטריצות}{matrix factorization} עדיפים על שיטות השכנים הקלאסיות. מודלים כאלה גם מאפשרים להכניס מידע נוסף, כמו משוב משתמע, השפעות של זמן ורמות ביטחון \cite{koren2009matrix}. הו, קורן ווולינסקי הראו כיצד להתאים מודל כזה במיוחד למשוב משתמע \cite{hu2008collaborative}.

האינטואיציה: במקום להשוות שורות ענקיות, המערכת ממפה כל משתמש וכל פריט ל\term{ייצוג וקטורי}{embedding}, כלומר לנקודה במרחב משותף. פריטים שאותם משתמשים צרכו מגיעים למקומות קרובים במרחב הזה \cite{google2026candidate}. הצירים של המרחב הזה נקראים \term{גורמים נסתרים}{latent factors}. כדי לדמיין את זה אפשר לחשוב על ציר כמו „כמה פופ משנות האלפיים". זו רק המחשה שלנו, אבל היא מסבירה למה נועה ויוני יוצאים קרובים, גם בלי שהמערכת יודעת שהם חברים.

\begin{inshort}
\begin{itemize}
  \item סינון שיתופי ממליץ לפי אנשים שהטעם שלהם דומה לשלך: מי שהסכימו בעבר כנראה יסכימו גם בעתיד.
  \item דמיון קוסינוס הוא מדד סטנדרטי לדמיון בין שני וקטורים. בדוגמה הבדיונית שלנו, יוני (0.95) קרוב לנועה יותר ממאיה (0.36).
  \item פירוק מטריצות מתאר משתמשים ושירים כנקודות במרחב משותף, ואפשר להתאים אותו למשוב משתמע.
\end{itemize}
\end{inshort}
